% APA 7th edition reference list (hand-formatted; keys match refs.bib).
% Page ranges use a hyphen to match the manuscript's dash convention (APA 7 itself specifies an en dash).
\begin{thebibliography}{21}
\setlength{\itemsep}{0pt}

\bibitem[{Avsec et~al.(2026)}]{avsec2026alphagenome}
Avsec, {\v Z}., Latysheva, N., Cheng, J., Novati, G., Taylor, K. R., Ward, T., Bycroft, C., Nicolaisen, L., Arvaniti, E., Pan, J., Thomas, R., Dutordoir, V., Perino, M., De, S., Karollus, A., Gayoso, A., Sargeant, T., Mottram, A., Wong, L. H., \ldots{} Kohli, P. (2026). Advancing regulatory variant effect prediction with AlphaGenome. \textit{Nature}, \textit{649}(8099), 1206-1218. \url{https://doi.org/10.1038/s41586-025-10014-0}

\bibitem[{Benegas et~al.(2023)}]{benegas2023gpn}
Benegas, G., Batra, S. S., \& Song, Y. S. (2023). DNA language models are powerful predictors of genome-wide variant effects. \textit{Proceedings of the National Academy of Sciences}, \textit{120}(44), Article e2311219120. \url{https://doi.org/10.1073/pnas.2311219120}

\bibitem[{Bredikhin et~al.(2026)}]{bredikhin2026alphagenomepytorch}
Bredikhin, D., Kjellberg, M., Zou, C., Buendia, A., Tu, X., \& Kundaje, A. (2026). \textit{alphagenome-pytorch: A PyTorch port of AlphaGenome} (Version 0.3.1) [Computer software]. GitHub. \url{https://github.com/genomicsxai/alphagenome-pytorch}

\bibitem[{Collins et~al.(2024)}]{collins2024tripod}
Collins, G. S., Moons, K. G. M., Dhiman, P., Riley, R. D., Beam, A. L., Van Calster, B., Ghassemi, M., Liu, X., Reitsma, J. B., van Smeden, M., Boulesteix, A.-L., Camaradou, J. C., Celi, L. A., Denaxas, S., Denniston, A. K., Glocker, B., Golub, R. M., Harvey, H., Heinze, G., \ldots{} Logullo, P. (2024). TRIPOD+AI statement: Updated guidance for reporting clinical prediction models that use regression or machine learning methods. \textit{BMJ}, \textit{385}, Article e078378. \url{https://doi.org/10.1136/bmj-2023-078378}

\bibitem[{Efron \& Tibshirani(1993)}]{efron1993bootstrap}
Efron, B., \& Tibshirani, R. J. (1993). \textit{An introduction to the bootstrap}. Chapman \& Hall.

\bibitem[{Feng et~al.(2025)}]{feng2025benchmark}
Feng, H., Wu, L., Zhao, B., Huff, C., Zhang, J., Wu, J., Lin, L., Wei, P., \& Wu, C. (2025). Benchmarking DNA foundation models for genomic and genetic tasks. \textit{Nature Communications}, \textit{16}, Article 10780. \url{https://doi.org/10.1038/s41467-025-65823-8}

\bibitem[{Geirhos et~al.(2020)}]{geirhos2020shortcut}
Geirhos, R., Jacobsen, J.-H., Michaelis, C., Zemel, R., Brendel, W., Bethge, M., \& Wichmann, F. A. (2020). Shortcut learning in deep neural networks. \textit{Nature Machine Intelligence}, \textit{2}, 665-673. \url{https://doi.org/10.1038/s42256-020-00257-z}

\bibitem[{Grimm et~al.(2015)}]{grimm2015circularity}
Grimm, D. G., Azencott, C.-A., Aicheler, F., Gieraths, U., MacArthur, D. G., Samocha, K. E., Cooper, D. N., Stenson, P. D., Daly, M. J., Smoller, J. W., Duncan, L. E., \& Borgwardt, K. M. (2015). The evaluation of tools used to predict the impact of missense variants is hindered by two types of circularity. \textit{Human Mutation}, \textit{36}(5), 513-523. \url{https://doi.org/10.1002/humu.22768}

\bibitem[{Hewitt \& Liang(2019)}]{hewitt2019control}
Hewitt, J., \& Liang, P. (2019). Designing and interpreting probes with control tasks. In \textit{Proceedings of the 2019 Conference on Empirical Methods in Natural Language Processing and the 9th International Joint Conference on Natural Language Processing (EMNLP-IJCNLP)} (pp. 2733-2743). Association for Computational Linguistics. \url{https://doi.org/10.18653/v1/D19-1275}

\bibitem[{Kao et~al.(2024)}]{kao2024lrb}
Kao, C.-H., Trop, E., Polen, M., Schiff, Y., de Almeida, B. P., Gokaslan, A., Pierrot, T., \& Kuleshov, V. (2024). Advancing DNA language models: The genomics long-range benchmark. In \textit{ICLR 2024 Workshop on Machine Learning for Genomics Explorations}. \url{https://mlanthology.org/iclrw/2024/kao2024iclrw-advancing/}

\bibitem[{Kapoor \& Narayanan(2023)}]{kapoor2023leakage}
Kapoor, S., \& Narayanan, A. (2023). Leakage and the reproducibility crisis in machine-learning-based science. \textit{Patterns}, \textit{4}(9), Article 100804. \url{https://doi.org/10.1016/j.patter.2023.100804}

\bibitem[{Karczewski et~al.(2020)}]{karczewski2020gnomad}
Karczewski, K. J., Francioli, L. C., Tiao, G., Cummings, B. B., Alf{\"o}ldi, J., Wang, Q., Collins, R. L., Laricchia, K. M., Ganna, A., Birnbaum, D. P., Gauthier, L. D., Brand, H., Solomonson, M., Watts, N. A., Rhodes, D., Singer-Berk, M., England, E. M., Seaby, E. G., Kosmicki, J. A., \ldots{} MacArthur, D. G. (2020). The mutational constraint spectrum quantified from variation in 141,456 humans. \textit{Nature}, \textit{581}(7809), 434-443. \url{https://doi.org/10.1038/s41586-020-2308-7}

\bibitem[{Landrum et~al.(2018)}]{landrum2018clinvar}
Landrum, M. J., Lee, J. M., Benson, M., Brown, G. R., Chao, C., Chitipiralla, S., Gu, B., Hart, J., Hoffman, D., Jang, W., Karapetyan, K., Katz, K., Liu, C., Maddipatla, Z., Malheiro, A., McDaniel, K., Ovetsky, M., Riley, G., Zhou, G., \ldots{} Maglott, D. R. (2018). ClinVar: Improving access to variant interpretations and supporting evidence. \textit{Nucleic Acids Research}, \textit{46}(D1), D1062-D1067. \url{https://doi.org/10.1093/nar/gkx1153}

\bibitem[{Ledneva et~al.(2026)}]{ledneva2026geneb}
Ledneva, D., Nuridinov, M., \& Kuznetsov, D. (2026). \textit{GENEB: Why genomic models are hard to compare} [Preprint]. arXiv. \url{https://arxiv.org/abs/2606.04525}

\bibitem[{Nguyen et~al.(2023)}]{nguyen2023hyenadna}
Nguyen, E., Poli, M., Faizi, M., Thomas, A., Birch-Sykes, C., Wornow, M., Patel, A., Rabideau, C., Massaroli, S., Bengio, Y., Ermon, S., Baccus, S. A., \& R{\'e}, C. (2023). HyenaDNA: Long-range genomic sequence modeling at single nucleotide resolution. In \textit{Advances in Neural Information Processing Systems} (Vol. 36). Curran Associates. \url{https://proceedings.neurips.cc/paper_files/paper/2023/hash/86ab6927ee4ae9bde4247793c46797c7-Abstract-Conference.html}

\bibitem[{Pedregosa et~al.(2011)}]{pedregosa2011sklearn}
Pedregosa, F., Varoquaux, G., Gramfort, A., Michel, V., Thirion, B., Grisel, O., Blondel, M., Prettenhofer, P., Weiss, R., Dubourg, V., Vanderplas, J., Passos, A., Cournapeau, D., Brucher, M., Perrot, M., \& Duchesnay, {\'E}. (2011). Scikit-learn: Machine learning in Python. \textit{Journal of Machine Learning Research}, \textit{12}, 2825-2830. \url{https://jmlr.org/papers/v12/pedregosa11a.html}

\bibitem[{Pejaver et~al.(2022)}]{pejaver2022calibration}
Pejaver, V., Byrne, A. B., Feng, B.-J., Pagel, K. A., Mooney, S. D., Karchin, R., O'Donnell-Luria, A., Harrison, S. M., Tavtigian, S. V., Greenblatt, M. S., Biesecker, L. G., Radivojac, P., \& Brenner, S. E. (2022). Calibration of computational tools for missense variant pathogenicity classification and ClinGen recommendations for PP3/BP4 criteria. \textit{American Journal of Human Genetics}, \textit{109}(12), 2163-2177. \url{https://doi.org/10.1016/j.ajhg.2022.10.013}

\bibitem[{Richards et~al.(2015)}]{richards2015acmg}
Richards, S., Aziz, N., Bale, S., Bick, D., Das, S., Gastier-Foster, J., Grody, W. W., Hegde, M., Lyon, E., Spector, E., Voelkerding, K., Rehm, H. L., \& ACMG Laboratory Quality Assurance Committee. (2015). Standards and guidelines for the interpretation of sequence variants: A joint consensus recommendation of the American College of Medical Genetics and Genomics and the Association for Molecular Pathology. \textit{Genetics in Medicine}, \textit{17}(5), 405-424. \url{https://doi.org/10.1038/gim.2015.30}

\bibitem[{Tang \& Koo(2025)}]{tang2025representational}
Tang, Z., \& Koo, P. K. (2025). Evaluating the representational power of pre-trained DNA language models for regulatory genomics. \textit{Genome Biology}. \url{https://doi.org/10.1186/s13059-025-03674-8}

\bibitem[{Tavtigian et~al.(2018)}]{tavtigian2018bayesian}
Tavtigian, S. V., Greenblatt, M. S., Harrison, S. M., Nussbaum, R. L., Prabhu, S. A., Boucher, K. M., \& Biesecker, L. G. (2018). Modeling the ACMG/AMP variant classification guidelines as a Bayesian classification framework. \textit{Genetics in Medicine}, \textit{20}(9), 1054-1060. \url{https://doi.org/10.1038/gim.2017.210}

\bibitem[{Zhou et~al.(2024)}]{zhou2024dnabert2}
Zhou, Z., Ji, Y., Li, W., Dutta, P., Davuluri, R. V., \& Liu, H. (2024). DNABERT-2: Efficient foundation model and benchmark for multi-species genomes. In \textit{The Twelfth International Conference on Learning Representations}. \url{https://iclr.cc/virtual/2024/poster/17823}

\end{thebibliography}
